\documentclass{beamer}

\usetheme{Madrid}

\usepackage[utf8]{inputenc}
\usepackage[brazil]{babel}
\usepackage{amsmath, amssymb}
\usepackage{xcolor}
\usepackage{tikz}
\usetikzlibrary{decorations.pathmorphing}

% Cor personalizada
\definecolor{mygreen}{RGB}{0,120,60}

% Aplicando cores e contrastes corretamente
\setbeamercolor{structure}{fg=mygreen}
\setbeamercolor{palette primary}{bg=mygreen, fg=white}
\setbeamercolor{palette secondary}{bg=mygreen, fg=white}
\setbeamercolor{palette tertiary}{bg=mygreen, fg=white}
\setbeamercolor{frametitle}{bg=mygreen, fg=white}
\setbeamercolor{title}{bg=mygreen, fg=white} 
\setbeamercolor{subtitle}{fg=white} 
\setbeamercolor{title in head/foot}{bg=mygreen, fg=white}

% Remove a barra inferior pesada do tema Madrid e deixa apenas o número do slide
\setbeamertemplate{footline}{%
  \hfill%
  \usebeamercolor[fg]{black}%
  \small\insertframenumber%
  \hspace*{3mm}\vskip3mm%
}
\setbeamertemplate{navigation symbols}{} 

% ==========================================
% ORGANIZAÇÃO DOS METADADOS
% ==========================================
\title{Processo de Gram-Schmidt}
\subtitle{Projeto Integrador na Extensão -- Ciência da Computação} 
\author{
João Daniel \\
Kelvin Krauss \\
Sofia Effting \\
Vinicius Vincent \\
}
\institute{Instituto Federal Catarinense - Campus Blumenau}
\date{29 de maio de 2026}

\begin{document}

% Slide 1 - Capa
\begin{frame}
\titlepage
\end{frame}

% --- PARTE 1: FUNDAMENTOS MATEMÁTICOS ---
\begin{frame}{Conceitos introdutórios}

% fundo cinza claro
\setbeamercolor{background canvas}{bg=gray!15}


\vspace{0.5cm}

% lista
\begin{itemize}
    \item Termos Utilizados
    \item O que é um Vetor?
    \item Definição de Ortogonalidade
    \item Operações Necessárias
    \item Regras para aplicação do GS
\end{itemize}

\end{frame}
%-----------------------------------
\begin{frame}{Termos Utilizados}

% fundo cinza
\setbeamercolor{background canvas}{bg=gray!15}

\begin{itemize}

    \item \textbf{Nomes dos vetores}
    \begin{itemize}
        \item $v$: vetores originais (entrada)
        \item $u$: vetores resultantes (saída)
    \end{itemize}

    \item \textbf{Projeção}
    \begin{itemize}
        \item Projeção é a componente de um vetor na direção de outro vetor.
    \end{itemize}
    
    \item \textbf{Número escalar}
    \begin{itemize}
        \item Um número escalar é um número que representa apenas uma quantidade (magnitude), sem direção.
    \end{itemize}

    \item \textbf{Espaço vetorial}
    \begin{itemize}
        \item É um conjunto de vetores, que podem ser somados entre si e multiplicados por números (chamados de escalares).
    \end{itemize}

\end{itemize}

\end{frame}
%-----------------------------------
\begin{frame}{O que é um Vetor?}

% fundo cinza claro
\setbeamercolor{background canvas}{bg=gray!15}

% conteúdo
\begin{itemize}
    \item Um vetor é uma grandeza que possui direção e magnitude, podendo ser representado graficamente como uma seta no plano.
    \item Ele pode existir em qualquer espaço, exemplos são:
    \begin{itemize}
        \item 2D: (x, y)
        \item 3D: (x, y, z)
    \end{itemize}
    \item Para dimensões superiores a 3 não é possível visualiza-los
\end{itemize}

\end{frame}
%-----------------------------------
\begin{frame}{Definição de Ortogonalidade}

\setbeamercolor{background canvas}{bg=gray!15}

\vspace{0.5cm}

\begin{itemize}
    \item Ortogonalidade generaliza o conceito de perpendicularidade
    \item Perpendicularidade entre vetores ocorre quando o angulo entre eles é reto, ou seja, 90° 
    \item Essa generalização amplia a ideia de perpendicularidade, transformando-a no conceito mais abrangente de ortogonalidade. 
    \item O qual determina que, em espaços vetoriais de qualquer dimensão, dois vetores $u$ e $v$ são ortogonais se o seu produto interno (escalar) for nulo, ou seja:
    \begin{itemize}
        \item $\langle u, v \rangle = 0$
    \end{itemize}
\end{itemize}

\vspace{0.5cm}

\end{frame}
%-----------------------------------
\begin{frame}{Definição de Ortogonalidade}

\setbeamercolor{background canvas}{bg=gray!15}

\begin{itemize}
    
    \item Geometricamente:
    \begin{itemize}
        \item Formam um ângulo de 90°
        \item Não possuem projeção um sobre o outro
    \end{itemize}
\end{itemize}

\vspace{0.8cm}

\begin{center}
\begin{tikzpicture}[scale=0.7]

% Eixos opcionais
\draw[->, gray] (-0.5,0) -- (4,0);
\draw[->, gray] (0,-0.5) -- (0,4);

% Vetor u
\draw[->, thick, blue] (0,0) -- (3,0)
node[below right] {$\vec{u}$};

% Vetor v
\draw[->, thick, red] (0,0) -- (0,3)
node[above left] {$\vec{v}$};

% Marca do ângulo reto
\draw (0.25,0) -- (0.25,0.25) -- (0,0.25);

% Texto 90 graus
\node at (0.7,0.7) {$90^\circ$};

\end{tikzpicture}
\end{center}

\end{frame}
%-----------------------------------
\begin{frame}{Operações Necessárias}

% fundo cinza
\setbeamercolor{background canvas}{bg=gray!15}

\vspace{0.5cm}

\begin{itemize}

\item \textbf{Produto escalar}
\begin{itemize}
    \item Gera um número escalar 
    \item Multiplica os elementos das coordenadas correspondentes e soma
    \item Exemplo: $\langle (2,3),(1,4) \rangle = 14$
\end{itemize}

\vspace{0.4cm}

\item \textbf{Subtração ou Soma de vetores}
\begin{itemize}
    \item Gera um vetor
    \item Ela faz a subtração ou soma elemento a elemento
    \item Exemplo: $(1,3) - (0,2) = (1,1)$, $(1,3) + (0,2) = (1,5)$
\end{itemize}

\vspace{0.4cm}

\item \textbf{Multiplicação de vetores por escalar}
\begin{itemize}
    \item Gera um vetor
    \item Multiplicação do escalar por todos os elementos do vetor
    \item Exemplo: $(1,3) * 2 = (2,6)$
\end{itemize}

\end{itemize}

\end{frame}
%-----------------------------------
\begin{frame}{Regras para aplicação do GS}

% fundo cinza
\setbeamercolor{background canvas}{bg=gray!15}

\vspace{0.5cm}

\begin{itemize}

\item \textbf{Vetores precisam ser independentes}
\begin{itemize}
    \item Um não pode ser formado a partir do outro
    \item Exemplo: $v_1 = (1,0)$ e $v_2 = (2,0)$
    \item Não pode ser usado pois $v_2 = 2v_1$
    \item Isso gera $u_2 = (0,0)$ no processo
\end{itemize}

\vspace{0.4cm}

\item \textbf{Vetores precisam ter a mesma dimensão}
\begin{itemize}
    \item Não é válido misturar dimensões diferentes
    \item Exemplo: $v_1 = (4,2)$ e $v_2 = (1,0,2)$
\end{itemize}

\vspace{0.4cm}

\item \textbf{A ordem dos vetores importa}
\begin{itemize}
    \item Resultados podem ser diferentes
    \item Mas todos continuam corretos
    %, pois os vetores finais permanecem ortogonais entre si, geram o mesmo subespaço vetorial e a base obtida ainda representa corretamente o espaço original
\end{itemize}

\end{itemize}

\end{frame}
%-----------------------------------
\begin{frame}{Produto Escalar}

\setbeamercolor{background canvas}{bg=gray!15}

\vspace{0.5cm}

\begin{itemize}
    \item O resultado é um número (valor escalar)
    \item Multiplicamos os elementos correspondentes dos vetores
    \item Em seguida, somamos os resultados dessas multiplicações
    \item Exemplo geral: $\langle (x_1, y_1), (x_2, y_2) \rangle = x_1x_2 + y_1y_2$
\end{itemize}

\vspace{0.5cm}

\[
\langle u, v \rangle = (x_1 \cdot x_2) + (y_1 \cdot y_2)
\]

\vspace{0.5cm}

\textbf{Exemplo:}
\[
\langle (2,3), (1,4) \rangle = (2 \cdot 1) + (3 \cdot 4) = 2 + 12 = 14
\]

\end{frame}
%-----------------------------------
\begin{frame}{Subtração de Vetores}

\setbeamercolor{background canvas}{bg=gray!15}

\vspace{0.5cm}

\begin{itemize}
    \item O resultado é um vetor
    \item Subtraímos os elementos correspondentes dos vetores
    \item Cada posição é calculada separadamente
    \item Exemplo geral: $(x_1, y_1) - (x_2, y_2) = (x_1 - x_2, y_1 - y_2)$
\end{itemize}

\vspace{0.5cm}

\[
(x_1, y_1) - (x_2, y_2) = (x_1 - x_2, y_1 - y_2)
\]

\vspace{0.5cm}

\textbf{Exemplo:}
\[
(1,3) - (0,2) = (1-0, 3-2) = (1,1)
\]

\end{frame}
%-----------------------------------
\begin{frame}{Soma de Vetores}

\setbeamercolor{background canvas}{bg=gray!15}

\vspace{0.5cm}

\begin{itemize}
    \item O resultado é um vetor
    \item Somamos os elementos correspondentes dos vetores
    \item Cada posição é calculada separadamente
    \item Exemplo geral: $(x_1, y_1) + (x_2, y_2) = (x_1 + x_2, y_1 + y_2)$
\end{itemize}

\vspace{0.5cm}

\[
(x_1, y_1) + (x_2, y_2) = (x_1 + x_2, y_1 + y_2)
\]

\vspace{0.5cm}

\textbf{Exemplo:}
\[
(1,3) + (0,2) = (1+0, 3+2) = (1,5)
\]

\end{frame}
%-----------------------------------
\begin{frame}{Multiplicação de escalares por vetor}

\setbeamercolor{background canvas}{bg=gray!15}

\vspace{0.5cm}

\begin{itemize}
    \item O resultado é um vetor
    \item Multiplicamos os elementos correspondentes dos vetores pelo escalar
    \item Cada posição é calculada separadamente
    \item Exemplo geral: $a * (x_1, x_2, \cdots,x_n)  = (ax_1, ax_2, \cdots,ax_n)$
\end{itemize}

\vspace{0.5cm}

\textbf{Exemplo:}
\[
4 * (0,2) = (4 \cdot 0, 4 \cdot 2) = (0,8)
\]

\end{frame}
%-----------------------------------
 % --- PARTE 1: FUNDAMENTOS MATEMÁTICOS ---

\begin{frame}{Introdução ao Processo de Gram-Schmidt}
\begin{itemize}
    \item \textbf{O que é?} Um algoritmo para ortogonalizar um conjunto de vetores em um espaço com produto interno.
    \item \textbf{Objetivo:} Transformar uma base qualquer $V$ em uma base ortogonal $U$.
    \item \textbf{Por que importa?} Bases ortogonais são mais simples.
\end{itemize}
\end{frame}
%----------------------------------
\begin{frame}{Fórmula da Projeção}
A projeção ortogonal de um vetor $v$ sobre um vetor $u$ é dada por:
\[
\text{proj}_{u}(v)= \frac {\langle v, u\rangle}{\langle u, u\rangle}u
\]
Para obter o vetor puramente ortogonal, subtraímos a projeção do vetor original:
\[
v_{\perp}=v-\text{proj}_{u}(v)
\]
Essa operação é o núcleo iterativo do processo de Gram-Schmidt.
\end{frame}
%-----------------------------------
\begin{frame}{Generalização Matemática}
Para um conjunto de vetores linearmente independentes $v_1, v_2, \dots, v_k$:
\[
u_1 = v_1
\]
\[
u_2 = v_2 - \text{proj}_{u_1}(v_2)
\]
\[
u_3 = v_3 - \text{proj}_{u_1}(v_3) - \text{proj}_{u_2}(v_3)
\]
\[
u_k = v_k - \sum_{j=1}^{k-1} \text{proj}_{u_j}(v_k)
\]
Cada passo remove todas as componentes paralelas aos vetores já processados.
\end{frame}
%-----------------------------------
\begin{frame}{Somatório (Sigma)}

\setbeamercolor{background canvas}{bg=gray!15}

\vspace{0.5cm}

\[
u_k = v_k - \sum_{j=1}^{k-1} \mathrm{proj}_{u_j}(v_k)
\]

\vspace{0.5cm}

\begin{itemize}
    \item $\sum$ é chamado de \textbf{somatório} (letra grega Sigma)
    \item Ele representa a soma de todas as operações que aconteceram do primeiro vetor ($j=1$) até o vetor imediatamente anterior ao atual ($k-1$).
\end{itemize}

\vspace{0.5cm}

\begin{itemize}
    \item $j=1$ → início da soma 
    \item $k-1$ → fim da soma 
\end{itemize}

\vspace{0.5cm}

\end{frame}
%-----------------------------------
\begin{frame}{Como o somatório funciona}

\setbeamercolor{background canvas}{bg=gray!15}

\vspace{0.5cm}

Expandindo o somatório:

\[
\sum_{j=1}^{k-1} \mathrm{proj}_{u_j}(v_k)
=
\mathrm{proj}_{u_1}(v_k)
+
\mathrm{proj}_{u_2}(v_k)
+
\cdots
+
\mathrm{proj}_{u_{k-1}}(v_k)
\]

\vspace{0.5cm}

\begin{itemize}
    \item Calculamos cada projeção separadamente
    \item Depois somamos todas elas
\end{itemize}

\vspace{0.5cm}

\begin{itemize}
    \item No final:
    \begin{itemize}
        \item subtraímos essa soma de $v_k$
        \item removemos todas as direções anteriores
    \end{itemize}
\end{itemize}

\end{frame}
%-----------------------------------
\begin{frame}{Construção da Fórmula da Projeção}
Queremos construir a projeção de $\vec{v}_2$ sobre $\vec{v}_1$:
\[
\mathrm{proj}_{\vec{v}_1}(\vec{v}_2)
\]

Vamos começar montando o círculo trigonométrico padrão: 
\end{frame}

%-----------------------------------
%========================================
% SLIDE 1 — Círculo trigonométrico
%========================================
\begin{frame}{Construção geométrica no círculo trigonométrico}

\centering

\begin{tikzpicture}[scale=1.2]

% círculo trigonométrico
\draw[thick] (0,0) circle (2);

% eixos
\draw[->] (-2.5,0) -- (2.5,0);
\draw[->] (0,-2.5) -- (0,2.5);

\end{tikzpicture}

\vspace{0.5cm}

\centering
1. Círculo trigonométrico

\end{frame}
%-----------------------------------
\begin{frame}{Construção da Fórmula da Projeção}

\setbeamercolor{background canvas}{bg=gray!15}

\vspace{0.7cm}

Depois vamos desenhar um vetor $\vec{v}_2$ partindo da origem até a borda do círculo.

\vspace{0.4cm}

Esse vetor forma um ângulo $\theta$ (theta) com o vetor $\vec{v}_1$, que está localizado no eixo $x$.

\vspace{0.4cm}
Por ultimo, iremos desenhar uma linha descendo reto da ponta do $\vec{v}_2$ até o chão, assim acabamos de montar um Triângulo Retângulo.
\vspace{0.4cm}

\end{frame}
%-----------------------------------
%========================================
% SLIDE 2 — Vetores com mesma origem
%========================================
\begin{frame}{Construção geométrica no círculo trigonométrico}

\centering

\begin{tikzpicture}[scale=1.2]

% círculo trigonométrico
\draw[thick] (0,0) circle (2);

% eixos
\draw[->] (-2.5,0) -- (2.5,0);
\draw[->] (0,-2.5) -- (0,2.5);

% pontos
\coordinate (O) at (0,0);
\coordinate (V1) at (2,0);
\coordinate (V2) at (1.2,1.6);

% vetor horizontal
\draw[->, thick, blue] (O) -- (V1)
node[below right] {$\vec{v}_1$};

% vetor inclinado
\draw[->, thick, red] (O) -- (V2)
node[above right] {$\vec{v}_2$};

\end{tikzpicture}

\vspace{0.5cm}

\centering
2. Vetores com mesma origem

\end{frame}
%========================================
% SLIDE 3 — Vetores + ângulo
%========================================
\begin{frame}{Construção geométrica no círculo trigonométrico}

\centering

\begin{tikzpicture}[scale=1.2]

% círculo trigonométrico
\draw[thick] (0,0) circle (2);

% eixos
\draw[->] (-2.5,0) -- (2.5,0);
\draw[->] (0,-2.5) -- (0,2.5);

% pontos
\coordinate (O) at (0,0);
\coordinate (V1) at (2,0);
\coordinate (V2) at (1.2,1.6);

% vetor horizontal
\draw[->, thick, blue] (O) -- (V1)
node[below right] {$\vec{v}_1$};

% vetor inclinado
\draw[->, thick, red] (O) -- (V2)
node[above right] {$\vec{v}_2$};

% arco do ângulo
\draw[thick]
(0.7,0)
arc[start angle=0,end angle=53,radius=0.7];

% theta
\node at (0.9,0.35) {$\theta$};

\end{tikzpicture}

\vspace{0.5cm}

\centering
3. Vetores + ângulo

\end{frame}


%========================================
% SLIDE 4 — Projeção horizontal
%========================================
\begin{frame}{Construção geométrica no círculo trigonométrico}

\centering

\begin{tikzpicture}[scale=1.2]

% círculo trigonométrico
\draw[thick] (0,0) circle (2);

% eixos
\draw[->] (-2.5,0) -- (2.5,0);
\draw[->] (0,-2.5) -- (0,2.5);

% pontos
\coordinate (O) at (0,0);
\coordinate (V1) at (2,0);
\coordinate (V2) at (1.2,1.6);
\coordinate (P) at (1.2,0);

% vetor horizontal
\draw[->, thick, blue] (O) -- (V1)
node[below right] {$\vec{v}_1$};

% vetor inclinado
\draw[->, thick, red] (O) -- (V2)
node[above right] {$\vec{v}_2$};

% projeção horizontal
\draw[dashed] (V2) -- (P);

% zig-zag verde na projeção
\draw[
    green!60!black,
    thick,
    decorate,
    decoration={
        zigzag,
        segment length=4pt,
        amplitude=1pt
    }
] (O) -- (P);

% triângulo
\draw[thin] (O) -- (V2) -- (P) -- cycle;

% arco do ângulo
\draw[thick]
(0.7,0)
arc[start angle=0,end angle=53,radius=0.7];

% theta
\node at (0.9,0.35) {$\theta$};

% ângulo reto
\draw (P) rectangle +(0.12,0.12);

\end{tikzpicture}

\vspace{0.5cm}

\centering
4. Projeção horizontal

\end{frame}

%-----------------------------------
\begin{frame}{Construção da Fórmula da Projeção}

\setbeamercolor{background canvas}{bg=gray!15}

\vspace{0.7cm}

Com o triangulo retângulo em mãos, nosso objetivo agora é descobrir o tamanho da "sombra" feita pelo vetor v2 no vetor v1.

\vspace{0.4cm}

Para descobrir a medida dessa sombra, utilizamos o cosseno de $\theta$, pois ele estabelece a relação exata entre o tamanho do vetor original e a base do triângulo.
\vspace{0.4cm}

\vspace{0.4cm}

\end{frame}

%-----------------------------------
\begin{frame}{Interpretação geométrica}

\setbeamercolor{background canvas}{bg=gray!15}

\vspace{0.4cm}

\begin{center}

\begin{tikzpicture}[scale=0.9]

% =====================================
% Fórmula do cosseno (lado esquerdo)
% =====================================

\node at (-4.2,1.5) {
\Large
$
\cos(\theta)
=
\frac{\text{cateto adjacente}}
{\text{hipotenusa}}
$
};

% =====================================
% Origem
% =====================================

\coordinate (O) at (0,0);

% Vetores
\coordinate (V1) at (4.8,0);
\coordinate (V2) at (2.7,3.1);

% Projeção
\coordinate (P) at (2.7,0);

% =====================================
% Vetor v1 em azul
% =====================================

\draw[->, thick, blue] (O) -- (V1)
node[below right] {$\vec{v}_1$};

% =====================================
% Vetor v2 em vermelho
% =====================================

\draw[->, thick, red] (O) -- (V2)
node[above right] {$\vec{v}_2$};

% =====================================
% zig-zag verde na projeção
% =====================================

\draw[
    green!60!black,
    thick,
    decorate,
    decoration={
        zigzag,
        segment length=4pt,
        amplitude=1pt
    }
] (O) -- (P);

% =====================================
% Altura da projeção
% =====================================

\draw[dashed] (V2) -- (P);

% =====================================
% Triângulo
% =====================================

\draw[thin] (O) -- (V2) -- (P) -- cycle;

% =====================================
% Ângulo reto
% =====================================

\draw (P) rectangle +(0.20,0.20);

% =====================================
% Arco do ângulo
% =====================================

\draw[thick]
(0.8,0)
arc[start angle=0,end angle=49,radius=0.8];

% =====================================
% Texto cos(theta)
% =====================================

\node at (1.38,0.45) {$\cos(\theta)$};

% =====================================
% Hipotenusa
% =====================================

\node[above left] at (1.55,1.75)
{\small Hipotenusa};

% =====================================
% Cateto adjacente
% =====================================

\node[below] at (1.35,-0.3)
{\small Cateto adjacente};

% =====================================
% Cateto oposto
% =====================================

\node[right] at (2.70,1.5)
{\small Cateto oposto};

\end{tikzpicture}

\end{center}

\end{frame}
%-----------------------------------

\begin{frame}{Cosseno no círculo trigonométrico}

\setbeamercolor{background canvas}{bg=gray!15}

\vspace{0.5cm}

Como $\vec{v}_2$ está no círculo trigonométrico:

\begin{itemize}
    \item A hipotenusa possui tamanho $1$
    \item Logo:
\end{itemize}

\vspace{0.5cm}

\[
\cos(\theta)
=
\frac{\text{cateto adjacente}}
{\text{hipotenusa}}
\quad
\Longrightarrow
\quad
\cos(\theta)
=
\frac{\text{cateto adjacente}}{1}
\]
\vspace{0.5cm}

Portanto:

\[
\boxed{
\text{cateto adjacente} = \cos(\theta)
}
\]

\vspace{0.4cm}

\begin{itemize}
    \item O cosseno representa a projeção horizontal do vetor
\end{itemize}

\end{frame}
%-----------------------------------
\begin{frame}{Vetores com tamanho qualquer}

\setbeamercolor{background canvas}{bg=gray!15}

\vspace{0.5cm}

\begin{itemize}

    \item Se o vetor $\vec{v}_2$ for maior ou menor que o círculo trigonométrico

    \item Basta multiplicar o cosseno pelo tamanho do vetor

\end{itemize}

\vspace{0.5cm}

\[
\boxed{
\text{cateto adjacente}
=
|\vec{v}_2| \cos(\theta)
}
\]

\vspace{0.4cm}

\end{frame}

%-----------------------------------
\begin{frame}{Usando produto escalar}

Sabemos que a definição geométrica do produto interno (ou escalar) entre dois vetores $\vec{v}_2$ e $\vec{v}_1$ é:

\[
 \langle \vec{v}_2 , \vec{v}_1 \rangle = |\vec{v}_2||\vec{v}_1|\cos\theta
\]
\end{frame}
%-----------------------------------
\begin{frame}{Isolando o termo}

\setbeamercolor{background canvas}{bg=gray!15}

\vspace{0.5cm}

Partimos da fórmula:

\[
\langle \vec{v}_2 , \vec{v}_1 \rangle = |\vec{v}_2||\vec{v}_1|\cos\theta
\]

\vspace{0.5cm}

\begin{itemize}
    \item Queremos isolar $|\vec{v}_2|\cos\theta$
    \item Para isso, dividimos ambos os lados por $|\vec{v}_1|$
\end{itemize}

\vspace{0.5cm}

\[
\frac{\langle \vec{v}_2 , \vec{v}_1 \rangle}{|\vec{v}_1|}
=
\frac{|\vec{v}_2||\vec{v}_1|\cos\theta}{|\vec{v}_1|}
\]


\begin{itemize}
    \item Simplificando $|\vec{v}_1|$:
\end{itemize}


\[
|\vec{v}_2|\cos\theta =
\frac{\langle \vec{v}_2 , \vec{v}_1 \rangle}{|\vec{v}_1|}
\]

\end{frame}
%-----------------------------------
\begin{frame}{Transformando a projeção em vetor}

\setbeamercolor{background canvas}{bg=gray!15}

\vspace{0.5cm}

Até agora, calculamos apenas o comprimento da projeção.

\begin{itemize}
    \item Mas queremos que a projeção seja um vetor
    \item Esse vetor deve apontar na mesma direção de $\vec{v}_1$
\end{itemize}


Para isso:

\begin{itemize}
    \item Pegamos o tamanho encontrado no passo anterior
    \item Multiplicamos por um vetor de tamanho $1$
    \item Esse vetor aponta na direção de $\vec{v}_1$
\end{itemize}

\vspace{0.6cm}

\end{frame}
%-----------------------------------
\begin{frame}{Transformando a projeção em vetor (Parte 2)}
Esse vetor especial é chamado de:

\[
\boxed{\text{Versor}} 
\]

\vspace{0.5cm}

O versor de $\vec{v}_1$ é dado por:

\[
\hat{v}_1 = \frac{\vec{v}_1}{|\vec{v}_1|}
\]

\end{frame}
%-----------------------------------
\begin{frame}{Construindo a projeção}

\setbeamercolor{background canvas}{bg=gray!15}

\vspace{0.5cm}

Agora juntamos:

\begin{itemize}
    \item Comprimento da projeção
    \item Direção da projeção
\end{itemize}

\vspace{0.5cm}

\[
\mathrm{proj}_{\vec{v}_1}(\vec{v}_2)
=
(|\vec{v}_2| \cos\theta)\hat{v}_1
\]

\end{frame}
%-----------------------------------
\begin{frame}{Interpretação}

\setbeamercolor{background canvas}{bg=gray!15}

\vspace{0.5cm}

\begin{itemize}
    \item $(|\vec{v}_2| \cos\theta)$ → tamanho da projeção
    \item $\hat{v}_1$ → direção da projeção
    \item Multiplicando os dois → obtemos um vetor
\end{itemize}

\end{frame}
%-----------------------------------
\begin{frame}{Substituição (parte 1)}

\setbeamercolor{background canvas}{bg=gray!15}

\vspace{0.5cm}

Começamos com:

\[
\mathrm{proj}_{\vec{v}_1}(\vec{v}_2)
=
(|\vec{v}_2|\cos\theta)\hat{v}_1
\]

\vspace{0.5cm}

Substituímos o vetor unitário:

\[
\hat{v}_1 = \frac{\vec{v}_1}{|\vec{v}_1|}
\]

\vspace{0.5cm}

\[
\mathrm{proj}_{\vec{v}_1}(\vec{v}_2)
=
(|\vec{v}_2|\cos\theta)
\left(\frac{\vec{v}_1}{|\vec{v}_1|}\right)
\]

\end{frame}
%-----------------------------------
\begin{frame}{Substituição (parte 2)}

\setbeamercolor{background canvas}{bg=gray!15}

\vspace{0.5cm}

Agora usamos:

\[
|\vec{v}_2|\cos\theta =
\frac{\langle \vec{v}_2 , \vec{v}_1 \rangle}{|\vec{v}_1|}
\]

\vspace{0.5cm}

Substituindo na expressão:

\[
\mathrm{proj}_{\vec{v}_1}(\vec{v}_2)
=
\left(\frac{\langle \vec{v}_2 , \vec{v}_1 \rangle}{|\vec{v}_1|}\right)
\left(\frac{\vec{v}_1}{|\vec{v}_1|}\right)
\]

\end{frame}
%-----------------------------------
\begin{frame}{Multiplicação das frações}

\setbeamercolor{background canvas}{bg=gray!15}

\vspace{0.5cm}

Partimos de:

\[
\left(\frac{\langle \vec{v}_2 , \vec{v}_1 \rangle}{|\vec{v}_1|}\right)
\left(\frac{\vec{v}_1}{|\vec{v}_1|}\right)
\]

\vspace{0.5cm}

\begin{itemize}
    \item Multiplicamos numerador com numerador
    \item Multiplicamos denominador com denominador
\end{itemize}

\vspace{0.5cm}

\[
=
\frac{(\langle \vec{v}_2 , \vec{v}_1 \rangle)\,\vec{v}_1}{|\vec{v}_1|\cdot|\vec{v}_1|}
\]

\vspace{0.5cm}

\begin{itemize}
    \item O produto escalar é um número (escalar)
    \item Ele multiplica o vetor $\vec{v}_1$
\end{itemize}

\end{frame}
%-----------------------------------
\begin{frame}{Simplificação}

\setbeamercolor{background canvas}{bg=gray!15}

\vspace{0.5cm}

Continuando:

\[
\frac{(\langle \vec{v}_2 , \vec{v}_1 \rangle)\,\vec{v}_1}{|\vec{v}_1|\cdot|\vec{v}_1|}
\]

\vspace{0.5cm}


Sabemos que:


\vspace{0.3cm}

\[
|\vec{v}_1|\cdot|\vec{v}_1| = \langle \vec{v}_1 , \vec{v}_1 \rangle
\]


\vspace{0.5cm}

Substituindo no denominador:

\[
=
\frac{(\langle \vec{v}_2 , \vec{v}_1 \rangle)\,\vec{v}_1}{\langle \vec{v}_1 , \vec{v}_1 \rangle}
\]

\end{frame}
%-----------------------------------
\begin{frame}{Reorganizando a expressão}

\setbeamercolor{background canvas}{bg=gray!15}

\vspace{0.5cm}

Temos:

\[
\frac{(\langle \vec{v}_2 , \vec{v}_1 \rangle)\,\vec{v}_1}{\langle \vec{v}_1 , \vec{v}_1 \rangle}
\]

\vspace{0.5cm}

\begin{itemize}
    \item $\langle \vec{v}_2 , \vec{v}_1 \rangle$ é um número (escalar)
    \item $\langle \vec{v}_1 , \vec{v}_1 \rangle$ também é um número
\end{itemize}

\vspace{0.5cm}

\begin{itemize}
    \item Ou seja, temos:
\end{itemize}

\vspace{0.3cm}

\[
\frac{\text{escalar} \cdot \vec{v}_1}{\text{escalar}}
\]

\end{frame}
%-----------------------------------
\begin{frame}{Regra algébrica}

\setbeamercolor{background canvas}{bg=gray!15}

\vspace{0.5cm}

Usamos a propriedade:

\[
\frac{a \cdot \vec{v}}{b} = \left(\frac{a}{b}\right)\vec{v}
\quad \text{(com $a,b$ escalares)}
\]

\vspace{0.5cm}

\begin{itemize}
    \item Escalares seguem as regras normais da divisão
    \item Podemos dividir $a$ por $b$ primeiro
    \item O resultado continua sendo um escalar
\end{itemize}

\vspace{0.5cm}

\begin{itemize}
    \item Esse escalar então multiplica o vetor
\end{itemize}

\vspace{0.5cm}

Aplicando:

\[
\text{proj}_{u}(v)= \frac {\langle v, u\rangle}{\langle u, u\rangle}u
=
\left(\frac{\langle \vec{v}_2 , \vec{v}_1 \rangle}{\langle \vec{v}_1, \vec{v}_1 \rangle}\right)\vec{v}_1
\]

\end{frame}
%-----------------------------------
\begin{frame}{Conclusão}

A fórmula da projeção vem diretamente de:
\begin{itemize}
    \item triângulo retângulo
    \item trigonometria
    \item produto escalar
\end{itemize}

\end{frame}
\begin{frame}{Exemplo 1: Duas Dimensões}
Considere a base em $\mathbb{R}^2$:
\[
v_1=(1,0), \quad v_2=(1,1)
\]
O primeiro vetor da nossa nova base ortogonal é simplesmente o primeiro vetor original:
\[
u_1 = v_1 = (1,0)
\]
\end{frame}
%-----------------------------------
\begin{frame}{Exemplo 1: O Produto Escalar}
Aplicando a fórmula da projeção para encontrar $u_2$:
\[
\text{proj}_{u_1}(v_2)
=
\frac{\langle v_2, u_1 \rangle}
{\langle u_1, u_1 \rangle}
u_1
\]
Calculando os produtos escalares:
\[
\langle v_2, u_1 \rangle
=
\langle (1,1) \cdot (1,0) \rangle
=
1\cdot1 + 1\cdot0
=
1
\]

\[
\langle u_1 , u_1 \rangle
=
\langle (1,0) \cdot (1,0) \rangle
=
1
\]
\end{frame}

\begin{frame}{Exemplo 1: Conclusão}
Substituindo na fórmula da projeção:
\[
\text{proj}_{u_1}(v_2) = \frac{1}{1}(1,0) = (1,0)
\]
Agora removemos a projeção do vetor original $v_2$:
\[
u_2 = (1,1) - (1,0)
\]
\[
u_2 = (0,1)
\]
Assim, a nova base ortogonal é formada por $u_1=(1,0)$ e $u_2=(0,1)$.
\end{frame}

\begin{frame}{Exemplo 2: Três Dimensões}
Considere os seguintes vetores em $\mathbb{R}^3$:
\[
v_1=(1,1,0)
\]
\[
v_2=(1,0,1)
\]
\[
v_3=(0,1,1)
\]
Definimos imediatamente o primeiro vetor:
\[
u_1=v_1=(1,1,0)
\]
\end{frame}
%-----------------------------------
%-----------------------------------
\begin{frame}{Exemplo 2: Calculando $u_2$}

Projeção de $v_2$ sobre $u_1$:

\[
\langle v_2, u_1 \rangle
=
\langle (1,0,1) \cdot (1,1,0) \rangle
=
1
\]

\[
\langle u_1 \cdot u_1 \rangle
=
\langle (1,1,0) \cdot (1,1,0) \rangle
=
2
\]

\[
\text{proj}_{u_1}(v_2)
=
\frac12(1,1,0)
=
\left(\frac12,\frac12,0\right)
\]

\[
u_2
=
(1,0,1)
-
\left(\frac12,\frac12,0\right)
=
\left(\frac12,-\frac12,1\right)
\]

\end{frame}

%-----------------------------------
\begin{frame}{Exemplo 2: Projeção de $v_3$ sobre $u_1$}

Para encontrar $u_3$, precisamos projetar $v_3$ sobre $u_1$ e $u_2$.

\vspace{0.5cm}

Primeiro, sobre $u_1$:

\[
\langle v_3, u_1 \rangle
=
\langle (0,1,1) \cdot (1,1,0) \rangle
=
1
\]

\[
\langle u_1, u_1 \rangle
=
\langle (1,1,0) \cdot (1,1,0) \rangle
=
2
\]

\[
\text{proj}_{u_1}(v_3)
=
\frac12(1,1,0)
\]

\end{frame}

%-----------------------------------
\begin{frame}{Exemplo 2: Projeção de $v_3$ sobre $u_2$}

Agora, sobre $u_2$:

\[
\langle v_3, u_2 \rangle
=
\left\langle
(0,1,1) \cdot
\left(
\frac12,
-\frac12,
1
\right)
\right\rangle
=
\frac12
\]

\[
\langle u_2, u_2 \rangle
=
\left\langle
\left(
\frac12,
-\frac12,
1
\right) \cdot
\left(
\frac12,
-\frac12,
1
\right)
\right\rangle
=
\frac14+\frac14+1
=
\frac32
\]

\[
\text{proj}_{u_2}(v_3)
=
\frac{\frac12}{\frac32}
\left(
\frac12,
-\frac12,
1
\right)
=
\frac13
\left(
\frac12,
-\frac12,
1
\right)
=
\left(
\frac16,
-\frac16,
\frac13
\right)
\]

\end{frame}
%-----------------------------------
\begin{frame}{Exemplo 2: Resultado Final de $u_3$}
Aplicando a fórmula completa para $u_3$:
\[
u_3 = v_3 - \text{proj}_{u_1}(v_3) - \text{proj}_{u_2}(v_3)
\]
\[
u_3 = (0,1,1) - \left(\frac12,\frac12,0\right) - \left(\frac16,-\frac16,\frac13\right)
\]
\[
u_3 = \left(-\frac23,\frac23,\frac23\right)
\]
A base $(u_1, u_2, u_3)$ agora é perfeitamente ortogonal.
\end{frame}
%-----------------------------------
\begin{frame}{Visualização Computacional: GeoGebra 2D}
Para comprovar a eficácia geométrica do algoritmo, desenvolvemos uma representação interativa no GeoGebra.
\vspace{0.3cm}
\begin{itemize}
    \item \textbf{O Setup:} Vetores originais $v_1$ e $v_2$ plotados no plano cartesiano.
    \item \textbf{A Dinâmica:} O software calcula a projeção.
    \item \textbf{O Resultado:} O vetor ortogonal $u_2$ reage, mantendo o ângulo exato de $90^\circ$ com $u_1$.
\end{itemize}
\vspace{0.5cm}
\end{frame}

%-----------------------------------
\begin{frame}{Visualização Computacional: GeoGebra 3D}
A visualização em três dimensões expõe a complexidade de isolar componentes espaciais.
\vspace{0.3cm}
\begin{itemize}
    \item \textbf{O Plano Base:} Os vetores ortogonalizados $u_1$ e $u_2$ formam um plano no espaço.
    \item \textbf{A Projeção Dupla:} Observamos o vetor $v_3$ sendo subtraído de suas componentes paralelas ao plano $(u_1, u_2)$.
    \item \textbf{O Vetor Normal:} O vetor $u_3$ resultante surge perfeitamente perpendicular à base formada.
\end{itemize}
\vspace{0.5cm}
\end{frame}

%-----------------------------------
\begin{frame}{Do Gráfico para a Prática}
O que o GeoGebra faz nos bastidores para renderizar essas retas ortogonais em tempo real?
\vspace{0.3cm}
\begin{itemize}
    \item O motor de renderização executa iterativamente as equações de projeção que acabamos de deduzir.
    \item Ele atualiza os cálculos a cada frame durante a interação do usuário.
    \item Isso demonstra que o Processo de Gram-Schmidt é o alicerce fundamental de motores geométricos, o que nos leva diretamente às aplicações práticas deste conceito.
\end{itemize}
\end{frame}

% --- PARTE 2: APLICAÇÕES PRÁTICAS ---

\begin{frame}{Aplicação 1: Orientação de Câmeras Virtuais (3D)}
Como motores gráficos (Unity, Unreal) mantêm a tela estável?
\vspace{0.3cm}
\begin{itemize}
    \item A câmera virtual é guiada por três vetores: \textit{Forward} (Frente), \textit{Up} (Cima) e \textit{Right} (Direita).
    \item Movimentos bruscos do mouse e acumulação de erros decimais fazem com que esses vetores percam seus ângulos exatos de $90^\circ$.
    \item Se a ortogonalidade for perdida, a imagem sofre "Shearing" (distorção). 
    \item O motor gráfico roda Gram-Schmidt silenciosamente a cada frame para re-ortogonalizar \textit{Up} e \textit{Right} em relação ao \textit{Forward}.
\end{itemize}
\end{frame}

%-----------------------------------
\begin{frame}{Aplicação 2: Ciência de Dados e Redundância}
Na Inteligência Artificial, vetores representam características de um dataset.
\vspace{0.3cm}
\begin{itemize}
    \item \textbf{O Problema:} Se usarmos "Número de Quartos" ($v_1$) e "Metragem Total" ($v_2$) para prever o preço de uma casa, os dados são redundantes (vetores quase paralelos).
    \item Redundância causa instabilidade e lentidão no treinamento de Redes Neurais.
    \item \textbf{A Solução:} Gram-Schmidt remove de $v_2$ tudo o que já é explicado por $v_1$.
    \item O novo vetor $u_2$ representa \textbf{apenas a informação nova e independente}, otimizando o modelo de Machine Learning.
\end{itemize}
\end{frame}

%-----------------------------------
\begin{frame}{Resumo do Projeto Integrador}
\begin{itemize}
    \item \textbf{Matemática:} O algoritmo limpa as dependências direcionais entre os dados de entrada, gerando bases perfeitamente independentes.
    \item \textbf{Computação:} É a ferramenta essencial para limpar dependências em grandes conjuntos de dados e estabilizar espaços virtuais em 3D.
    \item \textbf{Prática:} A união entre a base teórica e as soluções aplicadas formam a espinha dorsal de diversas tecnologias do nosso dia a dia.
\end{itemize}
\end{frame}

%-----------------------------------
\begin{frame}{Conclusão}
\begin{center}
\Large Obrigado pela atenção. \\
\vspace{0.5cm}
Abrimos agora para dúvidas.
\end{center}
\end{frame}

\end{document}
